\documentclass[11pt]{article}
\usepackage[a4paper,margin=25mm]{geometry}
\usepackage{amsmath,amssymb,amsthm,hyperref}
\newtheorem{theorem}{Theorem}
\newtheorem{lemma}{Lemma}
\newcommand{\PP}{\mathcal P}
\title{Positive completion preserving a fixed number of prefix moments}
\author{Complete proof; three independent mathematical reviews completed}
\date{30 September 2026}
\begin{document}\maketitle
Fix an integer $q\ge0$. Constants with subscript $q$ may depend on $q$,
which does not grow with the number $m\ge\max(1,q)$ of polynomial degrees.
Let $0=b_0<b_1<\cdots<b_K<b_{K+1}=1$ and
$\Delta=\min_i(b_{i+1}-b_i)$. Put
\[
 S_s(a)=K^{-1}\sum_{b_i>a}b_i^s,\qquad
 V_m=\operatorname{span}\{a^r:0\le r\le m\}
 +\operatorname{span}\{a^rS_s(a):r+s\le m\}.
\]
Let $P_q$ be orthogonal projection onto polynomials of degree at most $q$
separately on each interval $(b_i,b_{i+1})$, and put
\[
 \mathcal S_q=\operatorname{span}\{a^r:0\le r\le q\}
 +\operatorname{span}\{a^rS_s(a):0\le r\le q,\ r+s\le m\}.
\]
All norms and projections below use Lebesgue measure.

\begin{lemma}[A quotient estimate from matching derivative jumps]\label{quotient}
For every $F\in V_m$,
\[
 \operatorname{dist}_{L^2}(F,\mathcal S_q)
 \le C_q(m+1)^{2q+2}\Delta^{-q-1}\|(I-P_q)F\|_2.
\]
In particular $V_m\cap\operatorname{ran}P_q=\mathcal S_q$.
\end{lemma}
\begin{proof}
The jump of $F^{(\ell)}$ at $b_i$ is $-D_\ell(b_i)/K$, for a
polynomial $D_\ell$ of degree at most $m-\ell$.
For $0\le\ell\le q$ and $0\le t\le m-\ell$, define
\[
 \Phi_{\ell,t}(a)=\frac1{\ell!}
 \sum_{j=0}^{\ell}\binom\ell j(-1)^{\ell-j}
                      a^jS_{t+\ell-j}(a).
\]
Its jump polynomial at $b_i$ is
$-b_i^t(a-b_i)^\ell/(K\ell!)$.
Thus its derivatives have zero jump in every order except $\ell$,
where the jump is exactly $-b_i^t/K$; derivatives above $\ell$ vanish
inside every cell. Also $\Phi_{\ell,t}\in\mathcal S_q$.
Expanding each $D_\ell$ in powers and adding the corresponding
$\Phi_{\ell,t}$ constructs $G\in\mathcal S_q$ such that $F-G$ and its
first $q$ derivatives are continuous on $[0,1]$.

Subtract from $F-G$ its Taylor polynomial of degree $q$ at zero.
Repeated integration and Cauchy--Schwarz give
$\|F-G\|_2\le C_q\|F^{(q+1)}\|_2$, after absorbing that polynomial
into $G$. On each cell $I$, the fixed-order inverse Markov inequality
applied to $F-P_qF$ gives
\[
 \|F^{(q+1)}\|_{L^2(I)}
 \le C_q(m+1)^{2q+2}|I|^{-q-1}\|F-P_qF\|_{L^2(I)}.
\]
Summing squares proves the estimate. Its zero case gives the kernel
identity; the reverse inclusion is immediate.
\end{proof}

\begin{theorem}[Higher prefix completion]\label{main}
Suppose numbers $r_{s,i}$, $0\le s\le q$, $0\le i\le K+1$, satisfy
\[
 r_{s,0}=0,\qquad r_{s,K+1}=1/(s+1),
\]
and the compatibility identities
\begin{equation}\label{compat}
 K^{-1}\sum_{i=1}^K b_i^t r_{s,i}
 =\frac1{s+1}K^{-1}\sum_{i=1}^K b_i^{t+s+1}
 \qquad(0\le s\le q,\ s+t\le m).
\end{equation}
Put
$\delta=\max_{s,i}|r_{s,i}-b_i^{s+1}/(s+1)|$.
If
\[
 \delta\le c_q(m+1)^{-2q-3}\Delta^{2q+5/2},
\]
there is a density $f$, polynomial of degree at most $m$ on every cell,
such that $1/2\le f\le3/2$ and
\[
 \int_0^{b_i}a^s f(a)\,da=r_{s,i}\quad(0\le s\le q),
 \qquad \int fF=\int F\quad(F\in V_m).
\]
Consequently a variable of density $f$ and the equal-weight variable
on the nodes $b_i$ reproduce all ordered polynomial moments through
degree $m$ of a uniform variable and that same node variable.
If the node variable integrates a target density $w$ through degree
$m+1$, these are also the ordered moments of independent uniform and
$w$ variables.
\end{theorem}
\begin{proof}
First prescribe the desired moments through degree $q$ separately
on each cell. The difference from the uniform cell moment is at most
$2\delta$ in each raw monomial. On a cell $I=[a,a+\ell]$, expand in
the shifted Legendre basis of degrees $0,\ldots,q$, orthonormal for
the uniform probability measure on the cell.
Each of its polynomials has raw monomial coefficients bounded by
$C_q\ell^{-q}$, since $a\in[0,1]$ and $\ell\le1$.
The corresponding moment corrections therefore have size at most
$C_q\delta\ell^{-q}$. Their reproducing polynomial density, with the
normalizing factor $1/\ell$, is bounded by
$C_q\delta\ell^{-q-1}$. We obtain a piecewise polynomial $g$ of degree
$q$ with all the prescribed prefix moments and
\[
 \|g-1\|_\infty\le C_q\delta\Delta^{-q-1}.
\]
The endpoint conditions and \eqref{compat} say precisely that
$g-1$ is orthogonal to $\mathcal S_q$: for its step generators use
Fubini to get
$\int g a^sS_t=K^{-1}\sum_i b_i^t r_{s,i}$.

On $W=(I-P_q)V_m$, define
$\Lambda((I-P_q)F)=\int(g-1)F$.
Lemma~\ref{quotient} makes this well defined and bounds its norm by
$C_q(m+1)^{2q+2}\Delta^{-q-1}\|g-1\|_2$.
Let $h\in W$ be its Riesz representative. It has zero moments through
$q$ in every cell and is polynomial of degree at most $m$ there.
The cellwise polynomial evaluation estimate yields
\[
 \|h\|_\infty\le C(m+1)\Delta^{-1/2}\|h\|_2
 \le C_q(m+1)^{2q+3}\Delta^{-2q-5/2}\delta.
\]
Thus $f=g-h$ lies between $1/2$ and $3/2$ for sufficiently small $c_q$,
retains all prescribed cell moments, and integrates every $F\in V_m$
exactly as Lebesgue measure does.
For $r+s\le m$ this says
\[
 \mathbb E[A^rB^s\mathbf1_{A<B}]
 =\int f(a)a^rS_s(a)\,da
 =\frac1{r+1}K^{-1}\sum_i b_i^{r+s+1}.
\]
The ordinary and opposite-chamber moments follow as usual.
\end{proof}

\paragraph{Scope.}
For each fixed $q$, all loss factors are polynomial in $m$ and
$\Delta^{-1}$. The compatibility data still have to be constructed,
and polynomial common denominators have to be maintained across all
cells of any proposed induction. This lemma alone does not establish
an arbitrary fixed number of polynomial-size finite dice.
\end{document}
